% ============================================================
%  Section: Results
%  Figures:  Figure~\ref{fig:A}  -> H200 stats
%            Figure~\ref{fig:B}  -> L40S stats
%            Figure~\ref{fig:C}  -> RTX 4090 stats
%  All experiments: single-GPU, batch size 96, 200 epochs,
%  10 independent runs per configuration (except MONAI Base L40S, n=1).
% ============================================================

\section{Results}
\label{sec:results}

This section presents the quantitative evaluation of the six training
pipelines -- TensorFlow Base (\textbf{TF-Base}), TensorFlow Optimized
(\textbf{TF-Opt}), PyTorch Base (\textbf{PT-Base}), PyTorch Optimized
(\textbf{PT-Opt}), MONAI Base (\textbf{MN-Base}), and MONAI Optimized
(\textbf{MN-Opt}) -- across three GPU platforms: NVIDIA H200 SXM
(Figure~\ref{fig:A}), NVIDIA L40S (Figure~\ref{fig:B}), and NVIDIA
RTX~4090 (Figure~\ref{fig:C}).  All experiments were executed in
\textbf{single-GPU} mode, which has a critical implication: inter-GPU
communication fabric (NVLink, InfiniBand, PCIe switches) is entirely
inactive.  Therefore, every throughput difference observed between the
three platforms is attributable exclusively to on-chip resources --
memory bandwidth, Tensor Core density, and architectural features of
each GPU -- and not to any network or topology effect.

% -----------------------------------------------------------
\subsection{Throughput}
\label{subsec:throughput}

Table~\ref{tab:throughput_summary} and the throughput sub-panels in
Figures~\ref{fig:A}--\ref{fig:C} report mean training throughput in
images per second (img/s) over 10 independent runs.

\begin{table}[h]
  \centering
  \caption{Mean throughput (img/s) per framework, variant, and GPU platform.}
  \label{tab:throughput_summary}
  \begin{tabular}{lrrrrrr}
    \hline
    \textbf{GPU} & \textbf{TF-Base} & \textbf{TF-Opt} & \textbf{PT-Base} & \textbf{PT-Opt} & \textbf{MN-Base} & \textbf{MN-Opt} \\
    \hline
    H200    & 1\,053.8 & 3\,023.8 & 549.1 & 1\,924.0 & 1\,052.5 & 1\,803.2 \\
    L40S    &   476.1  & 1\,850.9 & 385.2 & 1\,099.2 &   840.3  & 1\,157.7 \\
    RTX 4090 &  518.8  & 1\,058.7 & 387.6 & 1\,097.5 &    ---   & 1\,177.5 \\
    \hline
  \end{tabular}
\end{table}

\paragraph{TensorFlow Opt.}
The application of DALI GPU decoding, FP16 mixed precision, and
XLA/JIT compilation (\texttt{--jit\_compile}) raised TF throughput from
476~img/s (TF-Base, L40S) to 1\,851~img/s (TF-Opt, L40S), a gain of
\textbf{+288.8\%} (Figure~\ref{fig:B}).  On the H200, the same
optimizations yield 3\,024~img/s starting from a Base of 1\,054~img/s,
a gain of \textbf{+186.9\%} (Figure~\ref{fig:A}).  Notably, the H200
achieved 3.024$\times$ the throughput of the L40S and 2.86$\times$ that
of the RTX~4090 under TF-Opt, even though all three results come from
single-GPU runs.  This disparity is explained by the H200's HBM3e
memory (4.8~TB/s bandwidth vs.\ 864~GB/s on the L40S), its higher
Tensor Core FP16 peak (${\sim}$989~TFLOPS vs.\ ${\sim}$366~TFLOPS on
the L40S), and the fact that the DALI + XLA pipeline, once data
movement from host to device is eliminated, becomes fully
compute-bound, allowing the H200's architectural advantages to
manifest completely.

The RTX~4090 under TF-Opt achieves only 1\,059~img/s (Figure~\ref{fig:C}),
a \textbf{+104.1\%} improvement over its Base of 519~img/s, but
\textbf{43\% less} than the L40S under the identical optimisation
stack.  Despite the RTX~4090 possessing marginally higher raw memory
bandwidth than the L40S (1\,008~GB/s vs.\ 864~GB/s), the L40S provides
${\sim}$11\% more FP16 Tensor Core compute (${\sim}$366 vs.\
${\sim}$330~TFLOPS) as a datacenter-grade chip with a wider 256-bit bus
that sustains higher throughput under sustained workloads without
consumer-grade power throttling.  The TF XLA kernel fusion is likewise
better exploited on professional Ada Lovelace GPUs, producing kernels
that fill the L40S's compute pipeline more efficiently than on the
RTX~4090.

\paragraph{PyTorch Opt.}
PT-Opt achieves a \textbf{+185.3\%} throughput gain on the L40S
(385~$\to$~1\,099~img/s) and a \textbf{+183.1\%} gain on the RTX~4090
(388~$\to$~1\,098~img/s) (Figures~\ref{fig:B}--\ref{fig:C}).  The two
values are \emph{practically identical} ($\Delta <$1\%), which contrasts
sharply with the TF-Opt case.  This parity occurs because
\texttt{torch.compile(mode=``reduce-overhead'')} with DALI and AMP
generates architecture-agnostic Triton kernels that overlap compute and
memory access at the op-level, leaving little room for the L40S's
datacenter features to differentiate.  On the H200, PT-Opt reaches
1\,924~img/s (+250.4\% over Base), once again due to the HBM3e
bandwidth feeding the PyTorch DALI pipeline faster than GDDR6/GDDR6X.

\paragraph{MONAI Opt.}
MONAI-Opt gains \textbf{+71.3\%} on the H200 (1\,053~$\to$~1\,803~img/s)
and \textbf{+37.7\%} on the L40S (840~$\to$~1\,158~img/s).  On the
RTX~4090, only the Opt variant was executed (1\,177~img/s), preventing
a percentage comparison.  The MONAI Opt improvements stem from combining
optional DALI preprocessing, cosine LR scheduling with warmup,
\texttt{channels\_last} (NHWC) tensor layout, and gradient
accumulation.  The NHWC format reorganises convolutional memory access
to align with Tensor Core paths, yielding fewer cache misses and
improved coalescence.  The H200 again benefits more (+71\%) than the
L40S (+38\%) because its higher memory bandwidth is better exploited
once the pipeline bottleneck shifts from CPU data loading to
GPU-resident computation.

% -----------------------------------------------------------
\subsection{Training Time}
\label{subsec:train_time}

\begin{table}[h]
  \centering
  \caption{Mean training time (seconds) per framework, variant, and GPU.}
  \label{tab:train_time_summary}
  \begin{tabular}{lrrrrrr}
    \hline
    \textbf{GPU} & \textbf{TF-Base} & \textbf{TF-Opt} & \textbf{PT-Base} & \textbf{PT-Opt} & \textbf{MN-Base} & \textbf{MN-Opt} \\
    \hline
    H200    & 2\,058 & 1\,024 & 5\,007 & 1\,488 & 2\,628 & 1\,385 \\
    L40S    & 4\,687 & 2\,243 & 6\,836 & 2\,629 & 3\,310 & 2\,242 \\
    RTX 4090 & 4\,376 & 3\,721 & 7\,032 & 2\,771 &   ---  & 2\,198 \\
    \hline
  \end{tabular}
\end{table}

\paragraph{PyTorch: largest absolute wall-time reduction.}
PT-Base is the slowest configuration on every platform.  The DALI +
AMP + \texttt{torch.compile} stack of PT-Opt cuts the training wall
time by \textbf{$-$70.3\%} on the H200 (5\,007~$\to$~1\,488~s),
\textbf{$-$61.5\%} on the L40S (6\,836~$\to$~2\,629~s), and
\textbf{$-$60.6\%} on the RTX~4090 (7\,032~$\to$~2\,771~s).  The
consistent 60--70\% reduction across all platforms indicates that
PT-Base's bottleneck is predominantly in the Python DataLoader and
CPU-side TFRecord decoding: once DALI moves decoding to the GPU and
\texttt{torch.compile} eliminates per-step Python dispatch overhead,
the GPU is kept busy throughout training regardless of which GPU is
used.  The extra 10\% gain seen on the H200 (+70\% vs.\ +60\%) is
attributable to the higher memory bandwidth allowing the DALI prefetch
buffer to be drained and refilled faster.

\paragraph{TensorFlow: aggressive pipeline, diminishing returns on RTX~4090.}
TF-Opt reduces training time by \textbf{$-$50.2\%} on the H200
(2\,058~$\to$~1\,024~s) and \textbf{$-$52.1\%} on the L40S
(4\,687~$\to$~2\,243~s), but only \textbf{$-$15.0\%} on the
RTX~4090 (4\,376~$\to$~3\,721~s).  This asymmetry mirrors the
throughput divergence discussed in Section~\ref{subsec:throughput}: the
XLA-compiled DALI pipeline does not extract the same latency
improvements from the consumer GPU architecture.  While the L40S
benefits from XLA kernel fusion with its larger shared-memory
configuration, the RTX~4090's TDP management and scheduler differences
limit sustained utilisation, preventing the same wall-time
compression.

\paragraph{MONAI: balanced reduction across platforms.}
MONAI-Opt reduces training time by \textbf{$-$47.3\%} on the H200
(2\,628~$\to$~1\,385~s) and \textbf{$-$32.3\%} on the L40S
(3\,310~$\to$~2\,242~s).  The more moderate gains compared to PT-Opt
and TF-Opt reflect MONAI's modular design, where DALI is optional and
not all transforms are offloaded to the GPU by default.  Even so, the
cosine scheduler and warmup allow convergence with fewer wasted
gradient steps, contributing to the time reduction beyond pure I/O
improvements.

% -----------------------------------------------------------
\subsection{Model Quality: AUC, Sensitivity, and Specificity}
\label{subsec:quality}

\begin{table}[h]
  \centering
  \caption{Mean AUC, Sensitivity (Sens.), and Specificity (Spec.) per variant and GPU.}
  \label{tab:quality_summary}
  \begin{tabular}{llrrr}
    \hline
    \textbf{Config} & \textbf{GPU} & \textbf{AUC} & \textbf{Sens.} & \textbf{Spec.} \\
    \hline
    TF-Base  & H200     & 0.985 & 0.891 & 0.981 \\
    TF-Base  & L40S     & 0.986 & 0.889 & 0.980 \\
    TF-Base  & RTX 4090 & 0.986 & 0.889 & 0.981 \\
    \hline
    TF-Opt   & H200     & 0.958 & 0.829 & 0.956 \\
    TF-Opt   & L40S     & 0.951 & 0.835 & 0.958 \\
    TF-Opt   & RTX 4090 & 0.958 & 0.820 & 0.970 \\
    \hline
    PT-Base  & H200     & 0.919 & 0.621 & 0.982 \\
    PT-Base  & L40S     & 0.918 & 0.641 & 0.980 \\
    PT-Base  & RTX 4090 & 0.919 & 0.603 & 0.986 \\
    \hline
    PT-Opt   & H200     & 0.985 & 0.887 & 0.981 \\
    PT-Opt   & L40S     & 0.985 & 0.888 & 0.980 \\
    PT-Opt   & RTX 4090 & 0.985 & 0.889 & 0.981 \\
    \hline
    MN-Base  & H200     & 0.971 & 0.856 & 0.972 \\
    MN-Base  & L40S     & 0.971 & 0.867 & 0.972 \\
    \hline
    MN-Opt   & H200     & 0.976 & 0.816 & 0.989 \\
    MN-Opt   & L40S     & 0.980 & 0.848 & 0.986 \\
    MN-Opt   & RTX 4090 & 0.976 & 0.801 & 0.989 \\
    \hline
  \end{tabular}
\end{table}

\paragraph{PyTorch: the only configuration where optimisation improves both speed and quality simultaneously.}
PT-Base produces an AUC of $0.918 \pm 0.004$ across all three GPUs,
with sensitivity of only $\approx$$0.62$, indicating the model
struggles with class imbalance when the data pipeline is slow and the
learning rate schedule is flat.  PT-Opt raises AUC to $0.985 \pm
0.001$ (\textbf{$+$0.067} absolute, all platforms), bringing
sensitivity to $\approx$$0.888$ and specificity to $\approx$$0.981$.
This simultaneous improvement in speed and quality is explained by the
synergy of three optimisations: (1)~DALI reduces starvation in the
DataLoader, yielding cleaner gradient estimates at each step; (2)~the
warmup + cosine LR scheduler together with mixup/label smoothing
regularise the larger effective batch produced by AMP without
introducing overfitting; and (3)~the EMA shadow model averages out
sharp loss spikes that occurred in PT-Base when the flat LR interacted
with imbalanced batches.  Because the backbone (\texttt{timm}
InceptionV3) and data pipeline are identical across GPUs, the AUC
reproducibility is near-perfect ($\sigma < 0.002$ on all three
platforms), demonstrating that PT-Opt's gains are architectural and
not dependent on GPU hardware.

\paragraph{TensorFlow: throughput-quality trade-off.}
TF-Base achieves the highest AUC ($0.985$--$0.986$) of any
configuration, consistent across all three GPUs (Figure~\ref{fig:A},
\ref{fig:B}, \ref{fig:C}).  This stability derives from the
\texttt{tf.data} pipeline's deterministic shuffling and the standard
Adam scheduler, which converge reliably with FP32 precision.
TF-Opt, however, reduces AUC by \textbf{$-$0.027} to $-$0.035
(H200: $-$0.027; L40S: $-$0.035; RTX~4090: $-$0.028).  The
degradation is caused by (1)~FP16 mixed precision introducing
occasional underflow/overflow in the cross-entropy loss for the
imbalanced diabetic-retinopathy classes; (2)~the aggressive
mixup/cutmix augmentation that blends images from both classes,
creating ambiguous intermediate examples that the model's final
sigmoid cannot resolve correctly without re-tuning; and (3)~the
progressive backbone unfreeze, which when combined with a short cosine
cycle, allows insufficient fine-tuning epochs for the deeper layers.
Sensitivity drops most noticeably ($0.889 \to 0.829$), confirming that
the positive class (referable retinopathy) is the primary victim of the
AUC loss.  Crucially, this trade-off is \emph{pipeline-driven} and not
GPU-dependent: the same AUC values are reproduced on all three
platforms.

\paragraph{MONAI: conservative but consistent gains.}
MN-Base delivers AUC $0.971$ on both the H200 and L40S.  MN-Opt
provides a modest but consistent improvement of \textbf{$+$0.005} to
$+$0.009 in AUC (H200: $0.976$; L40S: $0.980$; RTX~4090: $0.976$),
while sensitivity is reduced slightly ($\approx -0.04$ on the H200)
relative to MN-Base.  The AUC gain is attributable to the cosine LR
scheduler providing a larger effective learning rate during the early
epochs and the label smoothing ($\alpha = 0.02$) suppressing
overconfident predictions at the sigmoid output.  MONAI's modular
transform pipeline runs largely in CPU by default, which limits the
AUC benefit of DALI and explains the smaller quality delta compared to
PT-Opt, where the entire pipeline operates on GPU.

% -----------------------------------------------------------
\subsection{GPU Memory Consumption}
\label{subsec:memory}

\begin{table}[h]
  \centering
  \caption{Mean peak GPU memory (MB) per variant and GPU.}
  \label{tab:mem_summary}
  \begin{tabular}{lrrrrrr}
    \hline
    \textbf{GPU} & \textbf{TF-Base} & \textbf{TF-Opt} & \textbf{PT-Base} & \textbf{PT-Opt} & \textbf{MN-Base} & \textbf{MN-Opt} \\
    \hline
    H200    & 12\,216 & 1\,667 & 18\,967 & 8\,365 & 7\,282 & 6\,772 \\
    L40S    & 12\,358 & 4\,370 & 18\,911 & 18\,646 & 7\,400 & 6\,814 \\
    RTX 4090 & 10\,859 & 5\,097 & 18\,726 & 18\,600 &   ---  & 6\,814 \\
    \hline
  \end{tabular}
\end{table}

\paragraph{TensorFlow: the most dramatic memory reduction.}
TF-Opt reduces peak GPU memory by \textbf{$-$86.4\%} on the H200
(12\,216~MB~$\to$~1\,667~MB) and by \textbf{$-$64.6\%} on the L40S
(12\,358~MB~$\to$~4\,370~MB).  The discrepancy between the two GPUs
reflects a TensorFlow memory-caching behaviour: TF-Opt on the H200
takes advantage of the larger HBM3e capacity to release more
intermediate activations sooner (via XLA's rematerialisation pass),
achieving lower peak pressure, whereas the L40S's GDDR6 triggers more
conservative caching decisions.  The memory reduction itself is
structurally caused by three mechanisms: (1)~DALI processes images
entirely in GPU DRAM and does not stage FP32 copies per batch as
\texttt{tf.data} does; (2)~FP16 storage halves the activation tensor
footprint; and (3)~XLA graph fusion eliminates intermediate tensors
that are produced and consumed within a single fused kernel.

\paragraph{PyTorch: memory dominated by model size.}
PT-Base and PT-Opt both use the full InceptionV3 checkpoint and a
batch size of 96 images at $299{\times}299$, pushing peak memory to
$\approx$18.6~GB across all non-H200 platforms.  On the H200, PT-Opt
reduces memory to 8\,365~MB because \texttt{torch.compile} in
\texttt{reduce-overhead} mode performs graph-level fusion that
eliminates several large intermediate activation tensors, and the
gradient checkpointing introduced by the EMA mechanism recomputes some
activations during the backward pass instead of storing them.  On
GDDR6-based platforms (L40S and RTX~4090), this same fusion is still
applied but does not lower the peak below the minimum required for
AMP's combined FP16/FP32 parameter buffers at batch~96.

\paragraph{MONAI: stable memory footprint.}
MONAI-Base and MONAI-Opt maintain the smallest and most stable peak
memory of any framework ($7\,282$--$7\,400$~MB for Base;
$6\,772$--$6\,814$~MB for Opt), and the profile is identical across
platforms.  The slight \textbf{$-$7\%} reduction in Opt (6\,814 vs.\
7\,400~MB) is attributable to \texttt{channels\_last} tensors requiring
fewer padding bytes in convolutional layers, and to DALI not
materialising intermediate CPU-allocated NumPy arrays that MONAI's
default CPU transforms generate.

% -----------------------------------------------------------
\subsection{Cross-GPU Scaling: Why H200 Outperforms L40S and RTX~4090 in Single-GPU Mode}
\label{subsec:cross_gpu}

A recurring observation across all frameworks and variants is that the
H200 consistently delivers the highest throughput (Figure~\ref{fig:A}
vs.\ Figures~\ref{fig:B},~\ref{fig:C}).  Because all experiments use
exactly one GPU per run with no inter-GPU communication, this advantage
cannot be attributed to NVLink's GPU-to-GPU interconnect.  The
explanation is purely intra-chip:

\begin{enumerate}
  \item \textbf{HBM3e memory bandwidth.}
        The H200 SXM provides ${\sim}$4.8~TB/s of memory bandwidth vs.\
        $\approx$864~GB/s (L40S) and $\approx$1\,008~GB/s (RTX~4090).
        For memory-bound kernels -- including most normalisation,
        element-wise, and reduction operations -- the H200 can sustain
        computation at 4.7$\times$ the rate of the L40S and 4.8$\times$
        the rate of the RTX~4090.

  \item \textbf{Tensor Core compute.}
        The H200 delivers ${\sim}$989~TFLOPS at FP16, compared to
        ${\sim}$366~TFLOPS (L40S) and ${\sim}$330~TFLOPS (RTX~4090), a
        2.7$\times$ and 3.0$\times$ advantage in peak compute,
        respectively.  The DALI-based pipelines in TF-Opt and PT-Opt are
        the most compute-bound configurations, explaining why the H200
        advantage is larger under Opt variants than under Base variants.

  \item \textbf{Sustained throughput vs.\ peak throughput.}
        Consumer GPUs (RTX~4090) apply power throttling under sustained
        workloads to remain within their TDP envelope.  The L40S and
        H200, as datacenter-grade GPUs, sustain peak clocks for the full
        200-epoch run, which amplifies the throughput gap over long
        training jobs even if instantaneous FLOPS would suggest a smaller
        difference.

\end{enumerate}

The consequence is that TF-Opt scales \emph{super-linearly} relative
to L40S on the H200 (+186.9\% on H200 vs.\ +288.8\% on L40S, measured
from each GPU's own Base), while PT-Opt gains are consistent across
platforms (${\approx}$185--250\%).  The TF XLA/DALI pipeline is more
sensitive to memory bandwidth (super-linear scaling on HBM3e), whereas
PyTorch's Triton-generated kernels balance compute and memory access
more uniformly and therefore respond proportionally to raw compute
resources.

% -----------------------------------------------------------
\subsection{Summary of Baseline-to-Optimised Gains}
\label{subsec:summary_gains}

Table~\ref{tab:summary_gains} consolidates percentage changes from
Base to Opt for each framework and GPU, providing a compact reference
for the causal chains discussed in this section.

\begin{table}[h]
  \centering
  \caption{Summary of Base $\!\to\!$ Opt relative changes per GPU and framework.
           Memory and AUC are absolute changes; Throughput and Training Time are relative.}
  \label{tab:summary_gains}
  \resizebox{\columnwidth}{!}{%
  \begin{tabular}{llrrrr}
    \hline
    \textbf{Framework} & \textbf{GPU} &
    \textbf{Throughput} & \textbf{Train Time} & \textbf{AUC ($\Delta$)} & \textbf{Memory} \\
    \hline
    TensorFlow & H200     & +186.9\% & $-$50.2\% & $-$0.027 & $-$86.4\% \\
    TensorFlow & L40S     & +288.8\% & $-$52.1\% & $-$0.035 & $-$64.6\% \\
    TensorFlow & RTX 4090 & +104.1\% & $-$15.0\% & $-$0.028 & $-$53.1\% \\
    \hline
    PyTorch    & H200     & +250.4\% & $-$70.3\% & +0.067  & $-$55.9\% \\
    PyTorch    & L40S     & +185.3\% & $-$61.5\% & +0.068  & $-$1.4\%  \\
    PyTorch    & RTX 4090 & +183.1\% & $-$60.6\% & +0.066  & $-$0.7\%  \\
    \hline
    MONAI      & H200     & +71.3\%  & $-$47.3\% & +0.005  & $-$7.0\%  \\
    MONAI      & L40S     & +37.7\%  & $-$32.3\% & +0.009  & $-$7.9\%  \\
    \hline
  \end{tabular}}
\end{table}

The data in Table~\ref{tab:summary_gains} reveal three distinct
optimisation regimes:

\begin{itemize}
  \item \textbf{Regime 1 -- Throughput-first (TF-Opt):} gains exceed
        +100\% in throughput and $-$50\% in training time, but come at
        the cost of AUC ($-$0.027 to $-$0.035).  This regime is suitable
        when inference latency and GPU-hour cost dominate clinical
        deployment constraints and a small reduction in discriminative
        power is acceptable, provided the AUC remains above the
        clinical threshold.

  \item \textbf{Regime 2 -- Pareto-optimal (PT-Opt):} gains of
        +183\%--+250\% in throughput and $-$60\%--$-$70\% in training
        time are achieved with a simultaneous \emph{increase} in AUC
        (+0.066 to +0.068).  This is the only configuration that
        dominates its Base on all metrics and all platforms, making it
        the recommended default for resource-constrained single-GPU
        training of the HCPA retinopathy pipeline.

  \item \textbf{Regime 3 -- Balanced (MN-Opt):} modest throughput gains
        (+38\%--+71\%) with proportional training-time reductions and
        stable AUC (+0.005 to +0.009).  This regime is appropriate when
        MONAI-specific transforms or a clinical MONAI-based inference
        stack must be preserved, accepting lower absolute throughput in
        exchange for greater pipeline composability and lower memory
        pressure.
\end{itemize}
